\SetPageTheme{story}
\PageHeader{Lectia 2}{Ce se intampla cand impartim la 10}{L2 • P1}

\begin{missionbox}[De ce tocmai 10]
  Sistemul nostru de numeratie este zecimal. Asta inseamna ca fiecare cifra are legatura cu puterile lui 10. De aceea, impartirea la 10 produce un efect special: nu doar micsoreaza numarul, ci il reorganizeaza intr-un mod foarte util pentru algoritmi.

  Cand calculatorul trebuie sa citeasca un numar cifra cu cifra, impartirea la 10 devine una dintre cele mai importante operatii din toata lectia.
\end{missionbox}

\begin{conceptbox}[Regula de aur]
  Pentru orice numar natural \texttt{n}:
  \begin{itemize}
    \item \texttt{n \% 10} da ultima cifra;
    \item \texttt{n / 10} elimina ultima cifra.
  \end{itemize}
\end{conceptbox}

\begin{guidebox}[Primul pas]
  Pentru \texttt{n = 538}, ultima cifra este \hrulefill iar dupa eliminare ramane \hrulefill
\end{guidebox}

\newpage
\SetPageTheme{guided}
\PageHeader{Lectia 2}{Exemple clare pe numere}{L2 • P2}

\begin{examplebox}[Exemple]
  \begin{itemize}
    \item \texttt{1234 \% 10 = 4}
    \item \texttt{1234 / 10 = 123}
    \item \texttt{980 \% 10 = 0}
    \item \texttt{980 / 10 = 98}
    \item \texttt{7 \% 10 = 7}
    \item \texttt{7 / 10 = 0}
  \end{itemize}
\end{examplebox}

\begin{guidebox}[Exercitii ghidate]
  Completeaza:
  \begin{itemize}
    \item ultima cifra a lui 846 este \hrulefill
    \item dupa eliminarea ultimei cifre din 846 ramane \hrulefill
    \item ultima cifra a lui 1205 este \hrulefill
    \item dupa eliminarea ultimei cifre din 1205 ramane \hrulefill
  \end{itemize}
\end{guidebox}

\newpage
\SetPageTheme{guided}
\PageHeader{Lectia 2}{Observam tipare}{L2 • P3}

\begin{solobox}[Exercitii semi-ghidate]
  1. Completeaza tabelul:
  \begin{center}
  \begin{tabularx}{\linewidth}{|X|X|X|}
    \hline
    \textbf{n} & \textbf{n \% 10} & \textbf{n / 10} \\
    \hline
    451 & & \\
    \hline
    230 & & \\
    \hline
    9007 & & \\
    \hline
  \end{tabularx}
  \end{center}

  2. Scrie in cuvinte ce se intampla cu numarul atunci cand facem \texttt{/ 10}.
\end{solobox}

\begin{solobox}[Independent scurt]
  Inventeaza doua numere proprii si explica in doua etape cum afli ultima cifra si ce ramane dupa eliminare.
  \AnswerSpace{4}
\end{solobox}

\newpage
\SetPageTheme{challenge}
\PageHeader{Lectia 2}{Lucram fara sprijin}{L2 • P4}

\begin{solobox}[Fisa independenta]
  1. Calculeaza pentru fiecare numar ultima cifra si numarul ramas:
  \begin{itemize}
    \item 6729
    \item 4000
    \item 81
    \item 5
  \end{itemize}

  2. Explica de ce aceasta operatie este utila intr-un algoritm care vrea sa citeasca toate cifrele unui numar.
  \AnswerSpace{8}
\end{solobox}
